\documentclass[11pt]{article}

\usepackage[a4paper,margin=1in]{geometry}
\usepackage{amsmath,amssymb,amsthm}
\usepackage{booktabs}
\usepackage{hyperref}
\hypersetup{hidelinks}
\usepackage[T1]{fontenc}
\usepackage{lmodern}
\usepackage{xcolor}

\newtheorem{theorem}{Theorem}[section]
\newtheorem{lemma}[theorem]{Lemma}
\newtheorem{corollary}[theorem]{Corollary}
\newtheorem{proposition}[theorem]{Proposition}
\theoremstyle{definition}
\newtheorem{definition}[theorem]{Definition}
\newtheorem{remark}[theorem]{Remark}

\newcommand{\F}{\mathbb F}
\newcommand{\Zer}{\operatorname{Z}}
\newcommand{\etaext}{\eta}

\title{An Unbounded Additive Gap in a Minimum-Zero\\
Approach to Five-Cycle Double Covers}
\author{Atharva Vaidya\thanks{Draft prepared with substantial use of
AI systems.  See the disclosure in Section~\ref{sec:ai}.}}
\date{Draft of July 28, 2026}

\begin{document}
\maketitle

\begin{center}
\fbox{\parbox{0.92\textwidth}{
\textbf{Scope warning.}
This paper does \emph{not} prove or disprove the Five-Cycle Double Cover
Conjecture.  It disproves an unrestricted auxiliary assertion that was
proposed as one possible route to that conjecture.  Every graph constructed
here has an explicit five-cycle double cover.
}}
\end{center}

\begin{abstract}
For a cubic graph, a useful reformulation of the five-cycle double cover
problem asks for a matching that is simultaneously the exact zero set of an
\(\F_2^2\)-flow and the intersection of two binary cycles.  A natural
proof strategy is to minimize the number of zero edges first and then try
to extend a minimum exact-zero matching to such a certificate.  We give a
certified \(130\)-vertex counterexample to this strategy: its minimum
exact-zero matching size is \(5\), whereas every matching/flow/two-cycle
certificate has size at least \(6\), and one of size \(6\) exists.
We then prove an elementary factorization theorem for two-edge sums.
Iterating the sum gives connected simple bridgeless cubic graphs \(G_n\)
with
\[
 r_f(G_n)=r_M(G_n)=5\cdot2^n,\qquad
 \eta(G_n)=6\cdot2^n.
\]
Thus the additive extension gap is unbounded.  The infinite-family
argument is human-checkable; the finite base statement is backed by
explicit positive witnesses and LRAT certificates checked by two
independent proof checkers.
\end{abstract}

\section{Introduction and exact scope}

A binary cycle of a finite graph is an edge set having even incidence at
every vertex.  A standard five-cycle double cover (five-CDC) is an ordered
list of five binary cycles such that every edge belongs to exactly two
members of the list.  Empty members are allowed, so this convention also
expresses ``at most five''.  The Five-Cycle Double Cover Conjecture asserts
that every finite bridgeless graph has such a cover.  The conjecture is
due independently to Celmins and Preissmann; a current exposition of the
resolved unbounded-size cycle-double-cover problem explicitly records the
five-member strengthening as still open~\cite{Oum2026}.

This note studies only a proposed proof route for cubic graphs.  The route
starts with an \(\F_2^2\)-flow having as few zero edges as possible,
requires those zero edges to form a matching, and asks that the matching
extend to the intersection of two binary cycles.  The unrestricted
assertion that some cardinality-minimum exact-zero matching always extends
is false.  The purpose of this note is to make that negative result, and
its amplification under two-edge sums, precise and reproducible.

The result has two parts.
\begin{enumerate}
\item A finite, computer-assisted base theorem supplies a \(130\)-vertex
graph with extension gap \(6-5=1\).
\item A self-contained two-edge-sum theorem proves that both minima
factor exactly on identical rooted copies.  This turns the base graph into
an infinite family with unbounded additive gap.
\end{enumerate}
Neither part produces a five-CDC counterexample.  Indeed the size-six
certificate on the base graph, and its iterated gluing, give explicit
five-CDCs.

\paragraph{Publication status.}
This is a research draft, not a priority claim.  The literature comparison
in Section~\ref{sec:related} is sufficient to delimit the claim made here,
but a specialist review and independent human verification of the finite
certificate package are still required before submission.

\section{Definitions}

Throughout, graphs are finite and loopless.  The constructed graphs are
simple, although the definitions make sense for loopless multigraphs.
For an edge set \(D\), write \(\partial D\) for the set of vertices
incident with an odd number of edges of \(D\).  Thus \(D\) is a binary
cycle exactly when \(\partial D=\varnothing\).

\begin{definition}
An \(\F_2^2\)-flow on a graph \(G\) is a map
\(\phi:E(G)\to\F_2^2\) satisfying
\[
 \bigoplus_{e\ni v}\phi(e)=0
\]
at every vertex \(v\).  Its exact zero set is
\(\Zer(\phi)=\{e:\phi(e)=0\}\).
\end{definition}

\begin{definition}
For a cubic graph \(G\), define
\[
\begin{split}
r_f(G)&=\min_{\phi}|\Zer(\phi)|,\\
r_M(G)&=\min_{\substack{\phi\\ \Zer(\phi)\text{ a matching}}}
              |\Zer(\phi)|.
\end{split}
\]
If the second feasible set is empty, set \(r_M(G)=+\infty\).
\end{definition}

\begin{definition}
A matching/four-flow certificate on \(G\) is a quadruple
\((M,A,B,\phi)\) such that
\begin{enumerate}
\item \(M\) is a matching;
\item \(A\) and \(B\) are binary cycles;
\item \(A\cap B=M=\Zer(\phi)\); and
\item \(\phi\) is an \(\F_2^2\)-flow.
\end{enumerate}
Let
\[
 \etaext(G)=\min\{|M|:(M,A,B,\phi)
       \text{ is a matching/four-flow certificate}\},
\]
with value \(+\infty\) if there is no certificate.
\end{definition}

Plainly \(r_f(G)\le r_M(G)\le\etaext(G)\) whenever the quantities on the
right are finite.

\begin{proposition}[standard certificate equivalence {\cite[Corollary
0.6]{HoffmannOstenhof2013}}]
\label{prop:equivalence}
A finite loopless cubic multigraph has a standard five-CDC if and only if
it has a matching/four-flow certificate.
\end{proposition}

For completeness, the direction used below can be checked directly.
Given \((M,A,B,\phi)\), put \(D_0=A\triangle B\).  On \(G-M\), represent
the three nonzero values of \(\F_2^2\) by
\[
 s(1)=1100,\qquad s(2)=1010,\qquad s(3)=0110,
\]
and put \(k=1111\).  If \(\ell(y)\) denotes coordinate zero of \(s(y)\),
define
\[
 h(e)=1_{D_0}(e)+\ell(\phi(e)),\qquad
 d(e)=s(\phi(e))+h(e)k.
\]
Both summands defining \(h\) are binary cycles, so \(d\) is a
four-coordinate flow.  Each \(d(e)\) has weight two, and its coordinate
zero is \(1_{D_0}(e)\).  If \(D_0,D_1,D_2,D_3\) are the coordinate
supports, then
\[
 A,\ B,\ D_1,\ D_2,\ D_3
\]
double-cover every edge.  Conversely, pair-label the edges of a five-CDC;
at each cubic vertex the three labels form a triangle in \(K_5\).
Taking the intersection of two fixed coordinates gives the required
matching, and quotienting the remaining four even coordinates by
\(\langle1111\rangle\) gives the required \(\F_2^2\)-flow.

\section{The certified 130-vertex base graph}
\label{sec:base}

\subsection{Explicit construction}

Let \(Q\) be the \(60\)-vertex cubic graph whose graph6 encoding is
\begin{verbatim}
{??????O@?B?D?EGGSAK?H_?HG?Ao@A_??????????????C?G?C???M???B_???
[???????????????????O?C??C????B_????M?????[_???????????????????????
C??_???O?????B_?????B_?????@q?????????????????????????????C???G???
C???????M???????B_???????[_??????????????????????????????????O????
G???C????????B_????????M?????????[
\end{verbatim}
where the displayed line breaks are to be removed.  Mark the five edges
\[
 (22,26),\ (30,34),\ (38,42),\ (46,50),\ (54,58).
\]
Take two disjoint copies of \(Q\).  Subdivide the five marked edges in
each copy, respecting the displayed order, and join corresponding new
subdivision vertices.  Call the resulting graph \(G_0\).  It has
\(130\) vertices and \(195\) edges.  The five joining edges form a
matching.

The ancillary package
\begin{center}
\texttt{search/minimum-zero-exchange-countermodel-130v-20260727/}
\end{center}
contains the full edge list, graph6 encoding, positive witnesses, two CNF
instances, two LRAT proofs, a generator, and an independent verifier.

\begin{theorem}[certified base theorem]
\label{thm:base}
The graph \(G_0\) is simple, connected, bridgeless, and cubic, and
\[
 r_f(G_0)=r_M(G_0)=5,\qquad \etaext(G_0)=6.
\]
Moreover \(G_0\) has an explicit standard five-CDC.
\end{theorem}

\begin{proof}[Certificate proof]
The retained flow has exact zero set equal to the five joining edges.
Direct incidence checks verify flow conservation and verify that this zero
set is a matching.  Thus \(r_M(G_0)\le5\).

The first CNF asks for an arbitrary \(\F_2^2\)-flow with at most four
zero edges.  Its LRAT proof of unsatisfiability is accepted by both
\texttt{lrat-check} and CakeML \texttt{cake\_lpr}.  Hence
\(r_f(G_0)\ge5\).  Since \(r_f\le r_M\), the two inequalities prove
\[
 r_f(G_0)=r_M(G_0)=5.
\]

The second CNF is the complete matching/four-flow encoding with
\(|M|\le5\).  Its LRAT proof is accepted by the same two checkers, so
\(\etaext(G_0)\ge6\).  A retained size-six witness directly verifies
\[
 M=A\cap B=\Zer(\phi),
\]
the matching condition, the two binary-cycle conditions, and flow
conservation.  Therefore \(\etaext(G_0)=6\).  The package also stores the
five pair labels obtained by the reverse construction in
Proposition~\ref{prop:equivalence}; direct parity checking verifies that
the five coordinate supports double-cover every edge.

Finally, the verifier checks simplicity, cubicity, connectedness, and
bridgelessness directly from the reconstructed edge list.
\end{proof}

\begin{remark}
The computer-assisted part of Theorem~\ref{thm:base} is finite and
certificate-producing.  The positive assertions are checked directly.
The two negative assertions rely on LRAT proofs, not on trusting the SAT
solver that found them.  Appendix~\ref{sec:reproduce} gives replay
instructions and fixed hashes.
\end{remark}

\section{Exact factorization across a two-edge sum}

Let \(G_1,G_2\) be disjoint connected bridgeless cubic graphs with
distinguished edges
\[
 e_1=u_1v_1,\qquad e_2=u_2v_2.
\]
Delete \(e_1,e_2\) and add
\[
 f=u_1u_2,\qquad g=v_1v_2.
\]
Denote the resulting graph by \(G_1\#_2G_2\).  The crossed pairing has the
same properties.

\begin{lemma}[graph structure]
\label{lem:structure}
If the factors are connected, bridgeless, and cubic, then
\(G_1\#_2G_2\) is connected, bridgeless, and cubic.  If the factors are
simple, the sum is simple.
\end{lemma}

\begin{proof}
Cubicity, simplicity, and connectedness are immediate.  Each new edge
lies on the cycle consisting of \(f,g\) and a \(u_i\)-to-\(v_i\) path in
each \(G_i-e_i\); such paths exist because the root edges are not bridges.
For an old edge \(h\), choose a cycle of its factor containing \(h\).
If that cycle avoids the deleted root edge, it survives.  Otherwise,
replace the root edge on the cycle by the path through the other factor
and \(f,g\).  Hence every edge of the sum lies on a cycle, so the sum is
bridgeless.
\end{proof}

\begin{lemma}[cut parity]
\label{lem:cut}
Let \(D\) be a binary cycle in a graph and let
\(\{f,g\}\) be an edge cut.  Then \(D\) contains either both of \(f,g\)
or neither.
\end{lemma}

\begin{proof}
Sum the even-degree equations of \(D\) over one shore of the cut.
Internal incidences cancel, leaving
\(1_D(f)+1_D(g)=0\pmod2\).
\end{proof}

\begin{lemma}[flow restriction and zero count]
\label{lem:flow}
Let \(H=G_1\#_2G_2\), and let \(\phi\) be an \(\F_2^2\)-flow on \(H\).
Then \(\phi(f)=\phi(g)=a\) for some \(a\in\F_2^2\).  Restoring \(e_i\)
with value \(a\) gives a flow \(\phi_i\) on \(G_i\), and
\[
 |\Zer(\phi)|=|\Zer(\phi_1)|+|\Zer(\phi_2)|.
\]
If \(\Zer(\phi)\) is a matching, then each \(\Zer(\phi_i)\) is a matching.
\end{lemma}

\begin{proof}
Sum the flow equations over \(V(G_1)\).  Every internal edge appears
twice and cancels, leaving \(\phi(f)+\phi(g)=0\); in characteristic two
this says \(\phi(f)=\phi(g)=a\).  Assigning \(a\) to each restored root
edge makes the endpoint equations in each factor agree with the endpoint
equations in \(H\).

If \(a\ne0\), neither cut edge nor either restored root is zero.  If
\(a=0\), both cut edges and both restored roots are zero.  In either case,
the number of zeros in \(H\) is the sum of the numbers in the two
restored factors.

For the matching assertion, only the case \(a=0\) needs comment.  Since
the two zero cut edges meet the four root endpoints, no other zero edge of
\(\phi\) meets any of those endpoints.  Replacing the two cut edges by
one root edge in each factor therefore preserves the matching condition.
\end{proof}

\begin{lemma}[certificate restriction]
\label{lem:certificate}
Every matching/four-flow certificate on \(H=G_1\#_2G_2\) restricts to a
matching/four-flow certificate on each factor, and its matching size is
the sum of the two restricted matching sizes.
\end{lemma}

\begin{proof}
Let \((M,A,B,\phi)\) be a certificate.  By
Lemma~\ref{lem:cut}, each of \(A,B\) uses both cut edges or neither.
Restore \(e_i\) in each factor with the common membership bit for \(A\),
the common membership bit for \(B\), and the common flow value supplied
by Lemma~\ref{lem:flow}.  Binary-cycle parity and flow conservation are
restored at every root endpoint.  On a root edge the identities
\[
 1_M=1_A1_B,\qquad
 1_M=1_{\{\phi=0\}}
\]
hold because they held with the same states on \(f\) and \(g\).
Lemma~\ref{lem:flow} preserves the matching condition and proves the
additive count.
\end{proof}

\begin{theorem}[identical-rooted factorization]
\label{thm:factor}
Let \(G\) be a connected bridgeless cubic graph with a distinguished edge,
and let \(H=G\#_2G\) be formed from two identically rooted copies.  Then
\[
 r_f(H)=2r_f(G).
\]
If \(r_M(G)<\infty\), then
\[
 r_M(H)=2r_M(G).
\]
If \(\etaext(G)<\infty\), then
\[
 \etaext(H)=2\etaext(G).
\]
\end{theorem}

\begin{proof}
Lemma~\ref{lem:flow}, applied to an arbitrary flow on \(H\), gives
\[
 r_f(H)\ge2r_f(G).
\]
Take a minimum flow on \(G\), put identical copies on the two factors,
delete the two root edges, and give both new edges the common old root
value.  The endpoint equations hold and the zero-count identity gives
the reverse inequality.

If the zero set is required to be a matching, the same restriction and
gluing arguments preserve that condition by Lemma~\ref{lem:flow}.  This
proves the statement for \(r_M\).

Lemma~\ref{lem:certificate} gives
\(\etaext(H)\ge2\etaext(G)\).  Conversely, glue two identical copies of a
minimum certificate.  The two cycle-membership bits and the two flow bits
on the roots agree by construction, so replacing the roots with \(f,g\)
restores all parity equations and all certificate identities.  The
matching count doubles.
\end{proof}

\begin{remark}
For two nonisomorphic factors, the three restriction arguments still give
the corresponding additive lower bounds.  Equality can fail if minimum
objects in the two factors have incompatible states on their chosen root
edges.  Identical rooted copies avoid this compatibility issue.
\end{remark}

\section{The infinite family}

Set \(G_0\) equal to the graph of Theorem~\ref{thm:base}.  Having defined
a rooted \(G_n\), form \(G_{n+1}\) from two identically rooted copies by
the two-edge sum, and choose any surviving edge as the next distinguished
root.

\begin{corollary}[unbounded additive extension gap]
\label{cor:family}
For every \(n\ge0\), the graph \(G_n\) is finite, simple, connected,
bridgeless, and cubic, with
\[
\boxed{
 r_f(G_n)=r_M(G_n)=5\cdot2^n,\qquad
 \etaext(G_n)=6\cdot2^n.
}
\]
In particular,
\[
 \etaext(G_n)-r_M(G_n)=2^n
\]
is unbounded.  The graph \(G_n\) has \(130\cdot2^n\) vertices and
\(195\cdot2^n\) edges.
\end{corollary}

\begin{proof}
The structural assertions follow inductively from
Lemma~\ref{lem:structure}.  The three formulas follow by applying
Theorem~\ref{thm:factor} at every step, starting from
Theorem~\ref{thm:base}.  A two-edge sum deletes two edges and adds two, so
both the vertex and edge counts double.
\end{proof}

\begin{corollary}
Every \(G_n\) has a standard five-CDC.
\end{corollary}

\begin{proof}
The finite value of \(\etaext(G_n)\) supplies a matching/four-flow
certificate.  Apply Proposition~\ref{prop:equivalence}.
\end{proof}

\section{What the construction does not show}

The family is a counterexample family only to the unrestricted
minimum-zero extension strategy.  It is not a counterexample family to
the Five-Cycle Double Cover Conjecture.

For \(n\ge1\), the top-level composition exposes a cyclic two-edge cut.
The \(130\)-vertex base graph itself has girth \(4\) and cyclic
edge-connectivity \(3\).  Consequently this construction does not address
a version of the auxiliary assertion restricted to cyclically
four-edge-connected, high-girth graphs arising in standard
minimum-counterexample reductions.

The result establishes an unbounded \emph{additive} gap.  The ratio
\(\etaext(G_n)/r_M(G_n)=6/5\) is constant, so no claim of an unbounded
multiplicative gap is made.

The base theorem is computer-assisted.  The infinite factorization proof
does not depend on a solver, but its unconditional numerical starting
values depend on the correctness of the explicit finite witnesses, CNF
encodings, LRAT certificates, and proof-checker implementations.

\section{Related work and novelty boundary}
\label{sec:related}

The matching/four-flow characterization is not new.  Hoffmann-Ostenhof
proved that a cubic graph has a five-CDC if and only if it has a matching
\(M\), a nowhere-zero \(4\)-flow on \(G-M\), and two cycles whose exact
intersection is \(M\)~\cite[Corollary 0.6]{HoffmannOstenhof2013}.  Zhang
gave a broader characterization of five-cycle double covers through
covering subgraphs with nowhere-zero \(4\)-flows~\cite{Zhang1996}, and
Liu, Hao, Luo, and Zhang developed a \(4\)-flow and Catlin-reduction
approach for five-CDCs~\cite{LiuEtAl2023}.  Proposition~\ref{prop:equivalence}
is therefore included as the exact certificate language used by the
computation, not as a new existence characterization.

Flow resistance was formalized as an edge-uncolourability measure by
Fiol, Mazzuoccolo, and Steffen~\cite{FiolEtAl2018}.  Allie,
M\'a\v{c}ajov\'a, and \v{S}koviera constructed snarks with resistance
\(n\) and flow resistance
\(2n\)~\cite{AllieEtAl2022}, and later showed that the ratio between those
two established parameters can be arbitrarily large~\cite{AllieEtAl2024}.
Those results concern \(r_f\) versus ordinary resistance.  The parameter
\(\etaext\) here instead measures the least exact-zero matching that also
extends to the two-cycle certificate required for a five-CDC.

In the sources checked for this draft, we did not find the parameter
\(\etaext\), the assertion \(r_M=\etaext\), or the exact two-edge-sum
factorization in Theorem~\ref{thm:factor}.  The defensible contribution is
therefore narrow: a certificate-backed counterexample to that auxiliary
assertion and a human-checkable composition theorem making the additive
gap unbounded.  This is an apparent novelty statement, not a claim of
priority.  A domain expert should still check older flow, snark
superposition, and cycle-double-cover literature before submission.

\section{Reproduction and certificate boundary}
\label{sec:reproduce}

From the project root, run
\begin{verbatim}
cd search/minimum-zero-exchange-countermodel-130v-20260727
shasum -a 256 -c SHA256SUMS
python3 verify.py
\end{verbatim}
The verifier reconstructs the graph independently from the \(60\)-vertex
core and marked edges, checks graph structure and every positive witness,
regenerates both CNFs independently, compares them clause-for-clause with
the retained instances, and invokes both \texttt{lrat-check} and CakeML
\texttt{cake\_lpr} on both retained LRAT proofs.

The critical retained hashes are:
\begin{center}
\scriptsize
\begin{tabular}{ll}
\toprule
Artifact & SHA-256\\
\midrule
\texttt{graph.json} &
\texttt{a0cea6013ed45534441d072b8579792a643a6ad3b67f0948ec8a70eb0c5156b2}\\
\texttt{witness.json} &
\texttt{3dd958be4f8a6dbd88c2fbdca04e9904562468193db5e0c74f51f885adcfde79}\\
\texttt{flow-at-most-four-zero.cnf} &
\texttt{60c5c6a46aebf646d6d55c633052df56aa459aaf28a144b343a58bda0bf08397}\\
\texttt{flow-at-most-four-zero.lrat} &
\texttt{462544fce1c1d70797445136eaa9319aba159e42ae622ce9d229c662a909d626}\\
\texttt{extension-at-most-five.cnf} &
\texttt{df5f7e54866626c39c2f9c5a91b51a1ac4024174413707c1c700a4f319b6a6c8}\\
\texttt{extension-at-most-five.lrat} &
\texttt{e4a87af680d1a33d72f7e6825ab2f31f6a600eb87d200461776df51c5b5dd687}\\
\bottomrule
\end{tabular}
\end{center}

The first CNF has two flow-bit variables and one zero-indicator per edge.
It enforces both vertex-parity equations, the equivalence
\[
 z_e\longleftrightarrow
   (\neg p_e\land\neg q_e),
\]
and \(\sum_ez_e\le4\).

The second CNF has variables
\[
 m_e,\ a_e,\ b_e,\ p_e,\ q_e
\]
per edge.  It enforces
\[
 m_e\longleftrightarrow(a_e\land b_e),\qquad
 m_e\longleftrightarrow(\neg p_e\land\neg q_e),
\]
matching clauses
\(\neg m_e\lor\neg m_{e'}\) for incident pairs, four vertex-parity
equations for \(a,b,p,q\), and \(\sum_em_e\le5\).
Parity is expanded by its complete falsifying truth table; the cardinality
constraints use a sequential counter.  Thus satisfiability of the second
instance is exactly existence of a matching/four-flow certificate of size
at most five.

The additional script
\begin{center}
\texttt{preprint/minimum-zero-exchange-gap/audit\_two\_sum.py}
\end{center}
independently checks the positive base witnesses and, for every one of the
\(195\) possible root edges, constructs the two-copy sum and checks the
glued size-ten minimum-flow witness, the glued size-twelve certificate,
and the glued five-CDC.

\section{AI-use disclosure}
\label{sec:ai}

This research and this manuscript were produced with substantial
assistance from OpenAI Codex agents operating under human direction.  AI
systems participated in conjecture formulation, computational search,
construction discovery, SAT and certificate encoding, proof exploration,
independent-code drafting, audit suggestions, and substantial portions of
the prose and \LaTeX{} drafting.  The two-edge-sum proof was proposed and
checked in this AI-assisted process; it should not be represented as a
solely human derivation.

The use of AI does not replace mathematical verification.  The intended
trust boundary is explicit:
\begin{enumerate}
\item The factorization theorem and infinite-family corollary are
elementary arguments written in full in this paper and should be checked
line by line by human readers.
\item Positive finite claims reduce to public graph and witness data and
direct parity checks.
\item Negative finite claims reduce to fixed CNF files and fixed LRAT
proofs.  They are replayed by two proof checkers, one of which is the
CakeML verified checker used in the project package.
\item The manuscript has not yet received independent specialist peer
review.  Human verification, literature review, authorship review, and
venue-specific AI disclosure are required before submission.
\end{enumerate}
Any public version should preserve this disclosure and should list the
exact AI systems and versions if the submission venue requests them.

\begin{thebibliography}{9}

\bibitem{AllieEtAl2022}
I.~Allie, E.~M\'a\v{c}ajov\'a, and M.~\v{S}koviera,
\newblock Snarks with resistance \(n\) and flow resistance \(2n\),
\newblock \emph{Electron. J. Combin.} 29 (2022), no.~1, Paper P1.44,
\newblock \url{https://doi.org/10.37236/10633}.

\bibitem{AllieEtAl2024}
I.~Allie, E.~M\'a\v{c}ajov\'a, and M.~\v{S}koviera,
\newblock Flow resistance to resistance ratios in cubic graphs,
\newblock \emph{Discrete Appl. Math.} 342 (2024), 362--367,
\newblock \url{https://doi.org/10.1016/j.dam.2023.09.026}.

\bibitem{FiolEtAl2018}
M.~A. Fiol, G.~Mazzuoccolo, and E.~Steffen,
\newblock Measures of edge-uncolourability of cubic graphs,
\newblock \emph{Electron. J. Combin.} 25 (2018), no.~4, Paper P4.54.

\bibitem{HoffmannOstenhof2013}
A.~Hoffmann-Ostenhof,
\newblock A note on 5-cycle double covers,
\newblock \emph{Graphs Combin.} 29 (2013), 977--979,
\newblock \url{https://doi.org/10.1007/s00373-012-1169-8}.

\bibitem{LiuEtAl2023}
S.~Liu, R.-X.~Hao, R.~Luo, and C.-Q.~Zhang,
\newblock 5-cycle double covers, 4-flows, and Catlin reduction,
\newblock \emph{SIAM J. Discrete Math.} 37 (2023), 253--267,
\newblock \url{https://doi.org/10.1137/22M1472425}.

\bibitem{Oum2026}
S.-I.~Oum,
\newblock A proof of the cycle double cover conjecture by OpenAI:
an exposition,
\newblock arXiv:2607.16356v2 (2026),
\newblock \url{https://arxiv.org/abs/2607.16356}.

\bibitem{Zhang1996}
C.-Q.~Zhang,
\newblock Nowhere-zero 4-flows and cycle double covers,
\newblock \emph{Discrete Math.} 154 (1996), 245--253,
\newblock \url{https://doi.org/10.1016/0012-365X(95)00047-Z}.

\end{thebibliography}

\end{document}
